\section{Experiments}
\label{sec:experiments}

\paragraph{Setup.}
We use the official split files of the Grasp-Anything code: objects of \texttt{seen.obj} are cut 90\%/10\% into training and Seen test as in the authors' script, and \texttt{unseen.obj} is the Unseen test.
One prompt is one sample: 27.0k training, 3.0k Seen and 13.5k Unseen test prompts.
A prediction succeeds if its IoU with a ground-truth grasp exceeds 25\% and its angle differs by less than $30^{\circ}$~\cite{kumra2020antipodal}; H is the harmonic mean of Seen and Unseen.
Training follows the authors' code: Adam, learning rate $10^{-3}$, batch 8, 50 epochs of 1,000 batches; we use $\lambda = 0.1$ and $\tau = 0.07$ without tuning.
Both models are trained on two T4 GPUs with data parallelism (effective batch 16) under identical settings.

\paragraph{Dataset observations.}
On 1,000 random test objects with several parts, the part prompts of 70\% (Seen) and 83\% (Unseen) of the objects have near-identical grasp centres (within 5 pixels).
29\% and 39\% of the objects in the official Seen and Unseen files have no Grasp-Anything++ label.

\begin{table}[t]
  \centering
  \small
  \setlength{\tabcolsep}{4pt}
  \begin{tabular}{@{}lccc@{}}
    \toprule
    Method & Seen & Unseen & H \\
    \midrule
    CLIPORT~\cite{shridhar2021cliport}$^\dagger$ & 0.36 & 0.26 & 0.29 \\
    LGD~\cite{vuong2024language}$^\dagger$ & 0.48 & 0.42 & 0.45 \\
    MaskGrasp~\cite{vo2024maskgrasp}$^\dagger$ & 0.50 & 0.46 & 0.48 \\
    GraspMamba~\cite{nguyen2024graspmamba}$^\dagger$ & 0.63 & 0.41 & 0.52 \\
    \midrule
    GR-ConvNet + CLIP$^\dagger$ & 0.37 & 0.18 & 0.24 \\
    GR-ConvNet + CLIP, reproduced & 0.407 & 0.260 & 0.317 \\
    + region-text contrastive (ours) & \textbf{0.427} & \textbf{0.335} & \textbf{0.375} \\
    \bottomrule
  \end{tabular}
  \caption{Success rates on Grasp-Anything++. $^\dagger$ as reported by the authors (other data subsets); only the last two rows share our split and training setup. MaskGrasp's H is the harmonic mean of its reported Seen and Unseen.}
  \label{tab:main}
\end{table}

\begin{table}[t]
  \centering
  \small
  \setlength{\tabcolsep}{4pt}
  \begin{tabular}{@{}lcccc@{}}
    \toprule
    & \multicolumn{2}{c}{Baseline} & \multicolumn{2}{c}{Ours} \\
    Share of prompts & Seen & Unseen & Seen & Unseen \\
    \midrule
    Success & 40.7\% & 26.0\% & 42.7\% & 33.5\% \\
    Wrong object & 27.4\% & 35.5\% & 26.8\% & 33.7\% \\
    Right obj., wrong geometry & 20.8\% & 24.0\% & 22.2\% & 23.9\% \\
    No grasp predicted & 8.9\% & 11.5\% & 5.1\% & 4.9\% \\
    \bottomrule
  \end{tabular}
  \caption{Outcome of each test prompt. Failure types are judged by two vision-language models on 1,000 failures per model (87\% agreement; agreed cases only) and re-weighted to all failures; the small rest lands on the background.}
  \label{tab:errors}
\end{table}

\paragraph{Results.}
\Cref{tab:main} compares the methods and \cref{tab:errors} the failure types.
Trained under identical conditions, the contrastive loss improves the baseline by 2.0 points on Seen, 7.5 points on Unseen and 5.8 points of H.
Most of the gain comes from fewer missed predictions (no grasp on 11.5\% to 4.9\% of Unseen prompts), while wrong-object grasps drop only slightly (35.5\% to 33.7\%, \cref{tab:errors}), so grounding remains the main weakness.
